\documentclass[11pt,a4paper]{article}
\usepackage[margin=2.5cm]{geometry}
\usepackage{booktabs}
\usepackage{tabularx}
\usepackage{ragged2e}
\usepackage[T1]{fontenc}
\usepackage[usenames]{xcolor}
\usepackage[colorlinks=true,linkcolor=blue!60!black,urlcolor=blue!60!black]{hyperref}

\newcommand{\id}[1]{\texttt{#1}}
\newcolumntype{L}[1]{>{\RaggedRight\arraybackslash}p{#1}}

\renewcommand{\thesection}{S\arabic{section}}
\renewcommand{\thetable}{S\arabic{table}}
\setlength{\parskip}{4pt}

\begin{document}

\title{\textbf{Supplementary material for}\\[2pt]
``The Model Knows the Textbook but Not Your Class: Curriculum Knowledge
Graph Constraints for Slide Generation in Two Vocational-Education Courses''}
\author{Suxing Liu}
\date{}

\maketitle

\noindent This document holds eleven evidence and reference blocks that the main
text summarises: the ablation results (Section~\ref{sec:s1}), the mutation test
of the graph's own structural checks (Section~\ref{sec:s2}), the behavioural
probe of the generated code (Section~\ref{sec:s3}), the replication of the
claim-bearing conditions on a second vendor's model (Section~\ref{sec:s4}), the
concept-level paired analysis behind the per-baseline tests in \S5.4 of the main
text (Section~\ref{sec:s5}), the full qualitative readings behind the two cases
in \S5.8 of the main text (Section~\ref{sec:s6}), the verbatim prompt of one
session (Section~\ref{sec:s7}), one raw model response with the validator's view
of it (Section~\ref{sec:s8}), the seven measurement defects with symptom,
consequence and fix (Section~\ref{sec:s-defects}), the recomputation of the two
coverage conventions (Section~\ref{sec:s-coverage}), and the reproduction
commands (Section~\ref{sec:s11}). All numbers are regenerated by the scripts
named in each section from the archived raw outputs; nothing here is
hand-transcribed.

\section{Ablation of the contract}
\label{sec:s1}
% ==================================================================

Each ablation arm retains the complete contract and rule wording of \id{full}
and deletes exactly one data block, together with the rule text that would leak
the deleted information back in: \id{no\_order} removes the ordering block and
rule~3, \id{no\_misconception} removes the misconception block and rule~5, and
\id{no\_forbidden} removes the forbidden-term block and rule~6. The rule text
must be removed together with the data---deleting only the data while keeping a
rule that says ``you must follow the given order'' forces the model to guess an
order, and the ablation is no longer clean; a unit test guards this. All four
rows come from the ablation batch in \id{outputs/bench\_ablation/}---three
sessions $\times$ three resamplings, $n=9$ per arm---whose own nine \id{full}
re-runs anchor the table. This batch is \emph{separate} from the main
experiment, and comparisons are made within the table; on coverage and on
misconception coverage its \id{full} arm nonetheless lands on the same values as
the main batch (0.981 and 1.000), which is the agreement the main text relies on
when it places this arm beside the $2\times2$.

One arm in this table is a \emph{null manipulation} and is kept only for
completeness. Deleting the ordering block leaves the required-concept list, which
is itself emitted in topological order---the same sequence the deleted block
restates---so \id{no\_order} still receives the order. Its 0.000 violation rate is
therefore not evidence that an ordering constraint is unnecessary; it says only
that a second statement of the order can be removed without effect. A clean test
of the constraint would have to shuffle the concept list, which we did not do.

\begin{table}[htbp]
\centering
\caption{Ablations; each removes exactly one part of the graph. Mean (standard
deviation) over nine runs, all from the ablation batch.}
\label{tab:s-ablation}
\footnotesize
\setlength{\tabcolsep}{5pt}
\begin{tabularx}{\textwidth}{@{}L{3.2cm}cccc@{}}
\toprule
Condition & Coverage $\uparrow$ & Out-of-graph $\downarrow$ & Order viol. $\downarrow$
& Misconception $\uparrow$ \\
\midrule
\id{full} & 0.981 (0.056) & 0.000 (0.000) & 0.026 (0.063) & 1.000 (0.000) \\
$-$ ordering & \textbf{1.000} (0.000) & 0.000 (0.000) & 0.000 (0.000) & 1.000 (0.000) \\
$-$ misconceptions & 0.944 (0.083) & 0.000 (0.000) & 0.043 (0.065)
& \textbf{0.111} (0.333) \\
$-$ forbidden terms & \textbf{1.000} (0.000) & \textbf{0.083} (0.125) & 0.000 (0.000)
& 1.000 (0.000) \\
\bottomrule
\end{tabularx}
\end{table}

Regenerated by \id{python paper/tables.py} from \id{outputs/bench\_ablation/}.

% ==================================================================
\section{Mutation test of the graph's structural checks}
\label{sec:s2}
% ==================================================================

The shipped graph returns zero findings on all ten checks, which on its own is
weak evidence: a checker that never fires cannot be distinguished from one that
cannot fire. Each of the twelve codes the validator can emit is therefore
exercised by one minimal defect injected into a copy of the graph, and
Table~\ref{tab:s-checks} records whether the owning code responds and what
responds with it. All twelve codes fire on their own mutant. A single mutant
triggers a second code---an ungrounded skill also violates check~3---which is the
correct verdict rather than interference, so the checks are specific as well as
sensitive. Regenerated by \id{python paper/graph\_checks.py}.

\begin{table}[htbp]
\centering
\caption{Each of the twelve codes is exercised by one minimal defect injected
into a copy of the shipped graph. ``Also fires'' lists codes that respond in
addition to the intended one; a check that failed to fire here would be
decorative.}
\label{tab:s-checks}
\footnotesize
\begin{tabularx}{\textwidth}{@{}c l L{0.38\textwidth} L{0.19\textwidth}@{}}
\toprule
Chk & Code & Injected defect & Also fires \\
\midrule
1 & \texttt{DANGLING\_EDGE} & add relates\_to edge to undefined node C9999 & -- \\
\addlinespace[2pt]
2 & \texttt{PREREQ\_CYCLE} & reverse prerequisites\_for C1001$\to$C1003 & -- \\
\addlinespace[2pt]
3 & \texttt{NOT\_GROUNDED} & add Concept with no appears\_in edge & -- \\
\addlinespace[2pt]
4 & \texttt{MISSING\_EVIDENCE} & strip evidence from is\_a C1101$\to$C1001 & -- \\
\addlinespace[2pt]
5 & \texttt{UNIT\_CAPACITY} & add 6 filler concepts to U12 (5 present, cap 8) & -- \\
\addlinespace[2pt]
5 & \texttt{NO\_CORE\_CONCEPT} & downgrade the 3 mastery concepts of U12 & -- \\
\addlinespace[2pt]
6 & \texttt{DUPLICATE\_NAME} & rename C1204 to the name already used by C1001 & -- \\
\addlinespace[2pt]
7 & \texttt{FORBIDDEN\_REQUIRED\_OVERLAP} & forbid `circular queue' in U11, where it is a core concept & -- \\
\addlinespace[2pt]
8 & \texttt{FORBIDDEN\_SIBLING\_OVERLAP} & forbid `false overflow' in U11, where it is a sibling recap term & -- \\
\addlinespace[2pt]
9 & \texttt{ORPHAN\_MISCONCEPTION} & add Misconception with no addresses\_misconception in-edge & -- \\
\addlinespace[2pt]
10 & \texttt{IMPLEMENTS\_TARGET\_TYPE} & point implements at a Misconception node & -- \\
\addlinespace[2pt]
10 & \texttt{IMPLEMENTS\_SOURCE\_UNGROUNDED} & implements from a Skill with no appears\_in edge & \texttt{NOT\_GROUNDED} \\
\bottomrule
\end{tabularx}
\end{table}

Five unit tests in \id{tests/test\_graph.py} pin the four self-consistency
checks (7--10) to mutants of the shipped graph, so a regression in any of them
fails the suite.

% ==================================================================
\section{Behavioural probe of the generated code}
\label{sec:s3}
% ==================================================================

The probe (\id{paper/behavior\_check.py}) extracts every code block from the 54
main runs, strips top-level demo statements with an AST filter---otherwise a
demonstration loop leaves \id{front=2, rear=5} in the namespace and every later
probe fails for the wrong reason---executes the remainder in a fresh namespace,
and interrogates up to four behaviours: first-in-first-out order, the
empty-dequeue signal, the full signal, and survival of repeated enqueue. A block
qualifies only if it defines a module, a class, or a function whose first
parameter is a queue; fragments such as a bare \id{enqueue(self, x)} method body
are not testable. Class probes are scored separately for bounded and unbounded
implementations: an unbounded queue passes the full probe if 64 enqueues
complete without error, a bounded one if it eventually raises after at least one
successful enqueue. A probe counts as passed only when the code returns a
boundary signal (a sentinel, \id{False}, a raised exception, or a printed
message); a silent wrong answer fails.

\begin{table}[htbp]
\centering
\caption{Behavioural probe over all 117 code blocks from the 54 main runs.
``Testable'' is the share of blocks the harness can execute at all; the rate is
the share of testable blocks that satisfy every behaviour asked of them.}
\label{tab:s-probe}
\footnotesize
\begin{tabular}{@{}lcccc@{}}
\toprule
Condition & Blocks & Testable & Passed & Rate \\
\midrule
\id{full}      & 28 & 9  & 8  & 0.89 \\
\id{rag\_mc}   & 18 & 6  & 6  & 1.00 \\
\id{rag}       & 19 & 6  & 6  & 1.00 \\
\id{placebo}   & 18 & 3  & 3  & 1.00 \\
\id{cot}       & 18 & 7  & 7  & 1.00 \\
\id{zero\_shot}& 16 & 6  & 6  & 1.00 \\
\midrule
All            & 117 & 37 & 36 & 0.97 \\
\bottomrule
\end{tabular}
\end{table}

The single failure is a genuine scaffolding defect rather than a probe artefact:
in the U10 \id{full} run~3 deck, a program-style block writes
\id{enqueue(q, x)} against a capacity constant \id{CAP} that its own slide never
defines. This is precisely the ``slides are deliberately partial'' argument of
the main text: the excerpt assumes a surrounding program that a slide does not
carry.

As the main text argues, this is a bound and not a metric. The harness is
authored by the evaluators, so the evaluator supplies part of the answer it then
checks; and the qualifying denominator (0.32 of blocks, unevenly spread) is too
small and too condition-dependent to rank arms. \id{full} passing 8 of 9 against
28 of 28 for the baselines reads as a regression but is not evidence of one.

% ==================================================================
\section{Cross-model replication}
\label{sec:s4}
% ==================================================================

The same generator, prompts, sessions and validator, run on a second vendor's
model. This block exists to answer one question---does the contrast that carries
the paper (\id{full} against \id{rag\_mc}) keep its sign off the model it was
developed on---and it is not a comparison of model quality.

Two settings differ and are recorded per run in \id{outputs/bench\_kimi/results.csv}:
the backbone (\id{kimi-k3}) and the sampling temperature. Kimi K3 rejects any
temperature other than 1.0, so the replication runs at 1.0 against 0.3 for the
main table; the value is pinned in \id{ProviderPreset.fixed\_temperature} rather
than left to the caller's default, and a unit test holds it there. Within-model
contrasts are therefore unaffected, and we compare deltas across models, never
levels.

Runs: three sessions $\times$ two resamplings, six per condition, except
\id{zero\_shot} on Kimi ($n=5$): one cell could not be generated because the
vendor account reached its credit limit, and it is reported as missing rather
than refilled. \id{rag} and \id{cot} were not re-run on the second backbone; they
carry none of the claim and each extra condition costs about seven minutes at the
account's concurrency limit of one.

\begin{table}[htbp]
\centering
\caption{The same conditions under two backbones. Mean (standard deviation) over
runs; \id{rag} and \id{cot} were not re-run on the second backbone. DeepSeek V4 at
temperature 0.3, Kimi K3 at 1.0 (vendor-pinned), so levels are not comparable
across the two blocks---only deltas within a block are.}
\label{tab:s-cross}
\footnotesize
\setlength{\tabcolsep}{4pt}
\begin{tabular}{@{}l cccc c@{}}
\toprule
Condition & Coverage $\uparrow$ & Out-of-graph $\downarrow$ & Order viol. $\downarrow$ & Misconception $\uparrow$ & Pass \\
\midrule
\multicolumn{6}{@{}l}{\emph{DeepSeek V4}} \\
\id{full} & 0.981 (0.056) & 0.000 (0.000) & 0.020 (0.045) & 1.000 (0.000) & 67\% \\
\id{rag\_mc} & 0.698 (0.295) & 0.111 (0.220) & 0.007 (0.022) & 0.556 (0.527) & 22\% \\
\id{rag} & 0.659 (0.248) & 0.111 (0.220) & 0.030 (0.089) & 0.222 (0.441) & 0\% \\
\id{placebo} & 0.675 (0.266) & 0.167 (0.250) & 0.096 (0.131) & 0.333 (0.500) & 0\% \\
\id{cot} & 0.614 (0.185) & 0.019 (0.056) & 0.089 (0.183) & 0.389 (0.486) & 0\% \\
\id{zero\_shot} & 0.643 (0.248) & 0.019 (0.056) & 0.037 (0.089) & 0.222 (0.441) & 0\% \\
\midrule
\multicolumn{6}{@{}l}{\emph{Kimi K3}} \\
\id{full} & 1.000 (0.000) & 0.000 (0.000) & 0.016 (0.025) & 1.000 (0.000) & 67\% \\
\id{rag\_mc} & 0.734 (0.229) & 0.056 (0.136) & 0.033 (0.082) & 0.667 (0.516) & 17\% \\
\id{placebo} & 0.659 (0.188) & 0.056 (0.136) & 0.028 (0.068) & 0.333 (0.516) & 0\% \\
\id{zero\_shot} & 0.543 (0.270) & 0.000 (0.000) & 0.000 (0.000) & 0.400 (0.548) & 0\% \\
\bottomrule
\end{tabular}

\end{table}

Regenerated by \id{python paper/tables.py --cross outputs/bench\_kimi}.

\paragraph{What replicates.} On both metrics that carry the argument, and in the
same direction: misconception coverage $1.000 \to 0.667$ on Kimi against
$1.000 \to 0.556$ on DeepSeek, coverage $1.000 \to 0.734$ against
$1.000 \to 0.698$. The variance behaviour replicates more exactly than the means
do---\id{full} is 1.000 with zero standard deviation in every run under both
models, and \id{rag\_mc}'s standard deviation is 0.516 against 0.527.

\paragraph{What does not.} The gain shrinks on the second model ($-0.333$ against
$-0.444$ for misconceptions, $-0.266$ against $-0.283$ for coverage), and on U11
every condition scores 1.000 on Kimi, zero-shot included, so one of the three
sessions contributes no contrast at all. The pass-rate gap (67\% against 17\%)
rests on six runs per arm and carries no usable interval. Two backbones are two
draws from a large space: this rules out an artefact of one vendor's tuning and
supports no statement about model families.

% ==================================================================
\section{Concept-level paired analysis}
\label{sec:s5}
% ==================================================================

The main text reports coverage per deck, which leaves three clusters and a
clustered interval that includes zero. The contract, however, is written per
concept, and the same decks can be scored that way: every required concept is
either named in a deck or not, so each (session, resampling, concept) triple is a
paired binary observation shared by all conditions. Sixteen required concepts
across three sessions and three resamplings give $N=48$ pairs per baseline.

The matching is the shipped one (validator surface forms and graph aliases, body
pages only), so this section adds no new metric---only a different unit of
analysis. The test is the exact McNemar test on the discordant cells; the
per-unit counts are printed in the main text because the pairs are \emph{not}
independent (sixteen concepts nested in three sessions, and the three resamplings
of a cell repeat the same concept list), so the $p$ values describe the
consistency of the direction, not a population parameter.

\begin{table}[htbp]
\centering
\caption{Concept-level pairs against each baseline, $N=48$ per row. ``both'' =
covered by both conditions, ``\id{full} only'' / ``base only'' = the discordant
cells, RD = paired risk difference. Regenerated by
\id{python paper/concept\_pairs.py}.}
\label{tab:s-pairs}
\small
\begin{tabular}{@{}l rrrr rr@{}}
\toprule
Baseline & $N$ & both & \id{full} only & base only & RD & $p$ \\
\midrule
\id{rag\_mc} & 48 & 36 & 11 & 1 & $+0.208$ & $6.3\times10^{-3}$ \\
\id{rag} & 48 & 34 & 13 & 1 & $+0.250$ & $1.8\times10^{-3}$ \\
\id{placebo} & 48 & 35 & 12 & 1 & $+0.229$ & $3.4\times10^{-3}$ \\
\id{cot} & 48 & 30 & 17 & 1 & $+0.333$ & $1.4\times10^{-4}$ \\
\id{zero\_shot} & 48 & 33 & 14 & 1 & $+0.271$ & $9.8\times10^{-4}$ \\
\bottomrule
\end{tabular}

\end{table}

\begin{table}[htbp]
\centering
\caption{Leave-one-concept-out: the deck-level coverage gap against \id{rag\_mc}
is recomputed with each required concept removed in turn, per session. The range
is the answer to ``is one easy concept carrying the result''.}
\label{tab:s-loo}
\small
\begin{tabular}{@{}l ccc@{}}
\toprule
Session & gap, all concepts & range over drops & concepts \\
\midrule
U10 & $-0.056$ & $[-0.067, 0.000]$ & 6 \\
U11 & $+0.238$ & $[+0.111, +0.278]$ & 7 \\
U12 & $+0.667$ & $[+0.500, +1.000]$ & 3 \\
\bottomrule
\end{tabular}
\end{table}

\paragraph{Reading the pairs.} Two features matter. First, every baseline puts
more discordant cells on the side of \id{full} than on its own---11 to 1 against
\id{rag\_mc}, 17 to 1 against \id{cot}---and no baseline reverses the sign, so
the direction of the coverage effect does not depend on which control is chosen.
Second, no single concept carries it: dropping any one concept from the metric
moves a session's gap by at most 0.056, including U12, where three concepts alone
produce a gap between $+0.500$ and $+1.000$. What the pairs do \emph{not} fix is
U10: its discordant cells are 0 to 1 against \id{rag\_mc} and 3 to 1 against
\id{zero\_shot}, so the session-level sign can still flip there, exactly as the
clustered interval in the main text says.

% ===============================================================
\section{Qualitative readings of two cases}
\label{sec:s6}
% ===============================================================

\paragraph{A positive case: \id{full} on U12} U12 is the hardest session: it
requires choosing between a circular and a linked queue, a judgement task with no
single correct answer, and one required concept (``application scenarios of
queues'') is abstract and easily reduced to platitudes. All three \id{full} runs
reached coverage 1.000 and each produced a genuine comparison slide (predictable
capacity and frequent queries favour the circular queue; variable length and
frequent enqueue--dequeue favour the linked queue), handling the
breadth-first-search misconception on a dedicated counter-example slide. Coverage
for the five baselines on the same session is \textbf{0.333}: they cover only
``linked queue'' and dispose of selection and scenarios in one clause.

\paragraph{A negative case: \id{full} fails too} In the third resampling of U10,
coverage is 0.833, missing the concept ``queue'' itself. Reading slide by slide:
the title slide states that a queue is a linear list inserting at one end and
deleting at the other, but the first body slide is titled ``head, tail and the
sequential queue,'' and the deck never defines a queue separately again. The
model assumed students would remember the definition from one sentence on the
cover, and nothing in the following five slides reinforces it.

What kind of failure this is matters. Whether a definition on the title slide
counts as teaching is a pedagogical convention, and the omission disappears if the
title slide is admitted (coverage rises to 1.000); anyone who thinks a cover
definition suffices is entitled to read this run as a pass. What the metric
records is narrower: the definition is stated once, in a position the validator
does not treat as instruction, and never returned to. The validator flagged it,
which is the point of closing the loop---the method does not guarantee
correctness, it guarantees that \textbf{errors are not silently passed over}, and
it leaves the disposition to the teacher.

A second failure is convention-independent and therefore more instructive: the
second resampling of U10 contains two prerequisite inversions, the model teaching
the sequential queue first and returning to head and tail afterwards. On the
metrics this is a 0.133 violation rate; in the classroom it means students first
encounter the head and tail indices without knowing what they denote. The
topological order was in the prompt and the model violated it anyway, which is
why we insist on post-hoc validation rather than assuming that saying it implies
doing it.


\section{The prompt for one session}
\label{sec:s7}

The course is taught in Chinese, so the prompt is written in Chinese; that is a
property of the deployment, not a choice we made for the experiment. Verbatim
prompts for the \id{full} and \id{rag\_mc} conditions of session~U11 are
archived as \id{paper/prompt-U11-full.txt}; the rendering below keeps every
section header, every constraint and every hard rule, and condenses only the
per-item lists that the archive reproduces in full.

\begin{quote}\small
\textbf{Instruction.} Generate teaching slides according to the following
constraints.

\textbf{Session.} U11 Circular queue (7 pages).

\textbf{Teaching objective.} Understand how a circular queue reuses space
through modular arithmetic, distinguish the empty and full conditions
correctly, and write the implementation independently.

\textbf{Concepts that must be covered (none may be omitted).} Seven entries,
each with identifier, name, importance level, optional complexity and a
one-sentence definition, for example:
\id{C1104 Full condition (proficient)}: to separate ``full'' from ``empty'' one
storage cell is usually sacrificed, so the queue is full when
\id{(rear + 1) \% maxsize == front}.

\textbf{Explanation order (mandatory; this is the topological order induced by
the prerequisite relation).} The seven identifiers in sequence, with arrows.

\textbf{Misconceptions that must be corrected.} One entry: the erroneous claim
(the queue is full when \id{rear == front}), followed by the correction and the
consequence of confusing the two.

\textbf{Accompanying lab.} Implement a circular queue of capacity 8, run 20
enqueues and 15 dequeues, and verify modular reuse and the full condition.

\textbf{Evidence anchors (for teacher verification, not generated content).}
Six entries of the form ``\id{C1101: handout \S4.3: the circular queue} \id{exists
precisely to solve false overflow}''.

\textbf{Previously taught concepts that may be recapped in one sentence but not
expanded.} Ten sibling concept names, followed by the constraint that they must
not exceed 25\% of the text.

\textbf{Concepts that must not be introduced.} Five names (priority queue,
monotone queue, message queue, shortest path, topological sort).

\textbf{Output format.} A single JSON object with \id{unit\_id} and a
\id{slides} array; each slide carries \id{index}, \id{layout} (one of seven),
\id{title}, \id{bullets}, \id{code}, \id{notes}, \id{concept\_ids} and
\id{background}. No Markdown fence, no commentary.

\textbf{Hard requirements.} (1) Seven pages, first a title page and last a
summary. (2) Every concept in the mandatory list must appear in some page's
\id{concept\_ids}. (3) Page order must follow the given explanation order; no
prerequisite may be inverted. (4) At most four bullets per page, at most forty
characters each. (5) If misconceptions are listed, at least one page must have
layout \id{misconception}. (6) Do not introduce forbidden concepts; if one must
be mentioned, set \id{background} to true.
\end{quote}

The \id{rag\_mc} control replaces everything from \emph{Session} to
\emph{Forbidden concepts} with a flat list of the same definitions under the
heading \emph{Textbook excerpt}, and keeps the output format and hard rules
identical. That is the point of the control: the information is held constant
and only its structure changes.

\section{One raw response, and what the validator saw}
\label{sec:s8}

Responses are Chinese for the same reason the prompts are. The first two slides of
one \id{full} response (U11, run~1) are: index~1, layout \id{title}, ``Circular
queue'', three bullets naming the lecture, the objective and the lab, no code,
\id{concept\_ids} \id{[C1101]}; index~2, layout \id{concept}, ``The circular queue
and modular arithmetic'', four bullets on vacated cells, the logical ring,
\id{(index + 1) \% maxsize} and the $O(1)$ bound, code carrying the two
pointer-advance lines. The validator records coverage 1.000, prerequisite
violation 0.000, misconception coverage 1.000, out-of-graph 0.000, background
0.012 against a 0.25 ceiling, and editability 1.000, with no findings. The report
is addressed to concepts, not to a score: a failure names the uncovered concept
ids, the out-of-graph terms and the inverted pairs, so a low number can be traced
to a slide rather than only observed. The verbatim JSON and report are in the
archive under \id{outputs/main\_v1/}.

\section{The seven measurement defects}
\label{sec:s-defects}
% ==================================================================

The main text states the count, the common structure and the two readings a
reader should carry away; this section is the record. Seven defects were found
during the study and fixed before any number reported in the paper was
produced. None of them raised an error, and none was found by reflection: each
came from printing model outputs slide by slide and reading them against the
metric that scored them.

\begin{table}[htbp]
\centering
\caption{The seven measurement defects that silently favour the constraint-based
condition, and their fixes. Column ``Symptom'' names the surface observation;
``Consequence'' the wrong number it produced.}
\label{tab:s-defects}
\footnotesize
\setlength{\tabcolsep}{3pt}
\renewcommand{\arraystretch}{1.15}
\begin{tabularx}{\textwidth}{@{}cL{2.5cm}L{3.4cm}L{3.5cm}X@{}}
\toprule
\# & Defect & Symptom & Consequence & Fix \\
\midrule
1 & Verbatim coverage matching &
Graph says ``enqueue and dequeue operations'', baseline writes ``enqueue and
dequeue'' &
Baseline coverage crushed to 0.50 while \id{full} scores 1.00 by copying source
wording &
Add aliases to the graph; one shared vocabulary across conditions \\
2 & Order scored from self-reported IDs & Baselines never fill \id{concept\_ids} &
Empty sequence gives 0.00 violation; \textbf{an unanswered script scores full
marks} & Score from textual position instead \\
3 & Preview slides counted as body & Title and objectives slides list every
concept &
Ordering determined by list order rather than teaching order; three spurious
violations in U10 & Exclude title, objectives and summary \\
4 & Forward reference treated as teaching &
One slide remarks that modular arithmetic is key to implementing the circular
queue &
The concept is located three slides early; two spurious violations in U11 &
Locate by title naming first, fall back to first mention \\
5 & Cross-slide subtraction & The simple term is cancelled by a compound term on
another slide &
The concept is scored as never taught although the deck does teach it &
Compute net occurrences per line \\
6 & Misconception anchors too broad &
Unit-level \id{addresses\_\allowbreak misconception} used as a slide anchor &
Any slide with the same layout counts as corrected; the ablation still scores
1.00 & Restrict anchors to concept level; fix one mis-attached data edge \\
7 & Out-of-graph test too strict & Baselines invent their own ID scheme &
Out-of-graph rate inflated to 0.67 when the only difference is a namespace &
Count only non-existent IDs matching the graph's ID format \\
\bottomrule
\end{tabularx}
\end{table}

Three notes the table cannot carry. First, \textbf{the common structure across
all seven is that the representation the constrained method receives is exactly
the representation the metric uses}---this is the fairness problem of evaluating
a structured condition with unstructured metrics, and it is structural rather
than accidental, which is why it took seven separate sightings to see.

Second, every row is now pinned by a regression test in the file
\id{test\_measurement\_defects.py} under \id{tests/}: each test builds the deck
that used to fool the metric and asserts the value it should produce, and where
the wrong arithmetic can be written down, asserts that arithmetic too. The full
suite stands at 44 tests.

Third, defect~6 is the one that generalises beyond this study. It was not a bug
in metric code but \textbf{an error in the graph data}: one mis-attached
``related'' edge made any misconception-layout slide mentioning empty-queue
detection count as addressing a different misconception, so the ablation still
scored full marks and the conclusion about misconception injection nearly became
false. \textbf{Graph errors silently contaminate metrics}---a risk specific to
graph-driven methods, where the knowledge base is simultaneously the object of
study and part of the measuring instrument. We therefore recommend reporting how
the metrics behave \emph{on the baselines} and doing at least one manual pass
with outputs printed slide by slide.

% ==================================================================
\section{Sensitivity of the coverage convention}
\label{sec:s-coverage}
% ==================================================================

The coverage convention has two halves---body slides only, and matching on name
and aliases---and each is a decision rather than a fact about teaching, so each
was recomputed over all 54 main runs. Regenerated by
\id{python paper/coverage\_sensitivity.py}; the per-run values are in
\id{paper/coverage\_sensitivity.csv}.

\paragraph{Admitting the preview pages.} Keep the title slide, and coverage
moves in only two runs; the ranking of the six conditions is unchanged (Spearman
$\rho=1.00$; $\rho=0.94$ once objectives and summary are also admitted), and the
headline contrast against \id{rag\_mc} widens rather than shrinking
($+0.283\to+0.302$). The convention is therefore conservative with respect to the
claim: admitting the pages it excludes would strengthen the reported advantage,
not weaken it.

\paragraph{Name-only matching.} The aliases decision points the same way, and
more sharply. Recomputing coverage with name-only matching drops baseline
coverage to 0.325--0.389 while \id{full} falls only to 0.963, because \id{full}
was handed the canonical strings and copied them (thirty-one of the 54 runs move;
the largest single drop, from 1.000 to 0.143). Alias matching therefore
\emph{narrows} the reported gap: the conservative number for \id{full} against
\id{rag\_mc} is $+0.283$, and the gap a stricter matcher would produce is
$+0.574$. This is the opposite of a defect that manufactures the
result---it is a correction that runs against the hypothesis, which is why we
keep it and why the reported advantage should be read as a lower bound rather
than a best case.

% ==================================================================
\section{The repair loop}
\label{sec:s-repair}
% ==================================================================

The validator emits findings, and in the deployed workflow a teacher acts on
them; in the main experiment nothing did. One third of the \id{full} runs
contain at least one error-level finding, so we ran a repair loop on the stored
runs of two cells: the DeepSeek grid of the main results table (nine runs, three
failures) and the \id{agnes-2.5-flash} cell of the sweep (row~3; six runs, two
failures)---the cell where the misconception slot's advantage disappears, and
therefore the harder place to ask.

The policy is fixed and deliberately plain: for a run with an error-level
finding, the validator's error findings are appended to the same condition
prompt and the whole deck is regenerated, up to three attempts; first-pass runs
are untouched, so the repaired count shares its denominator with the first-pass
count. The appended text names the uncovered concepts and the inverted pairs and
adds nothing else---for \id{full} those names are already in the contract, so a
finding restates a requirement instead of supplying one. The same loop on
\id{zero\_shot} would hand over the syllabus that condition exists to lack, which
is why it is not run there: whether a repair loop is fair is a property of the
condition, not of the loop.

\paragraph{On the main grid the loop closes the gap.} All three failures were
repaired in a single round each: $6/9 \to 9/9$, mean coverage $0.981 \to 1.000$,
three model calls in total. The repairs are substantive: a missing concept was
given the page it lacked (in U10/r2 the prerequisite \id{C1002} existed only
inside a later summary bullet), and an inverted pair was fixed by splitting a
page that had merged two concepts (U11/r3), so the metrics are computed on more
content than before, not less.

\paragraph{On the third backbone it does not.} The same policy on the
\id{agnes-2.5-flash} cell moved $4/6$ to $5/6$: the missing-concept failure was
repaired in one round, the ordering failure was not repaired in three, and its
violation rate went $0.048 \to 0.048 \to 0.381$, worse on the third attempt than
on the first.

\paragraph{Reading the split.} \textbf{The loop closes every failure on the
backbone that fills a required slot and one of two on the backbone that does
not}; the same division appears for the slot itself in the second course. A
repair loop is not a separate capability that rescues a weak condition---it
inherits the contract's dependence on the backbone. Two cautions. Three successes
out of three is a count, not a rate: its exact 95\% lower bound is
$0.05^{1/3} \approx 0.37$, so the loop is shown to work and not shown to be
reliable. And three attempts on one deck produced three different decks with no
convergence, so a repair budget cannot rest on the expectation of convergence; it
has to be set by cost. The main table therefore keeps 67\% as a first-pass
figure, and $9/9$ is reported as the outcome of a bounded, costed repair on nine
decks---not as the pass rate of the method.

% ==================================================================
\section{Reproducibility}
\label{sec:s12}

All runs were issued against the vendor's hosted API
(\url{https://api.deepseek.com}) with model id \id{deepseek-v4-flash},
temperature 0.3, the provider's default thinking mode, and no fixed seed; the
model therefore returns different slide text on repeated calls, which is why each
cell is resampled three times rather than run once. Slides are emitted as JSON
and rendered with \id{python-pptx}; no image generation model is involved, which
is what makes \id{editability} 1.000 in every condition.


Code, graph, experimental protocol and all raw outputs are available in the
\id{kg-slidegen} repository. The accompanying protocol document records how every
number in this paper was produced, and each experimental cell retains three
artefacts: the presentation file, the raw model output and a per-item validation
report. The teacher ratings are released as the file
\id{paper/human\_eval.csv}, one row per deck--rater pair, and
\id{paper/human\_eval.py} regenerates the blinded ratings table in the main text
from it, together with the ICC values and their bootstrap intervals.

% ===============================================================

\end{document}
